\chapter{Diagnóstico del modelo}\label{cap:diagnostico}

\begin{entregable}{(d): diagnóstico}
Se contrastan las hipótesis del modelo lineal sobre los residuos del modelo final: linealidad,
independencia de los errores, homocedasticidad, normalidad y ausencia de multicolinealidad, y se
identifican las observaciones influyentes. La \cref{tab:diagnostico} resume los contrastes; para
cada violación se indica cómo la absorbe la estrategia de inferencia.
\end{entregable}

Bajo las hipótesis clásicas, $\mathbf y=X\boldsymbol\beta+\boldsymbol\varepsilon$ con
$\E(\boldsymbol\varepsilon\mid X)=\mathbf 0$ y $\Var(\boldsymbol\varepsilon\mid X)=\sigma^2 I$, el
estimador MCO $\hat{\boldsymbol\beta}=(X^{\top}X)^{-1}X^{\top}\mathbf y$ es el de menor varianza
entre los lineales insesgados (teorema de Gauss-Markov) y, si además los errores son normales, los
contrastes $t$ y $F$ son exactos. Conviene separar qué hipótesis afecta a qué:
\begin{itemize}
  \item la \textbf{linealidad} (media condicional bien especificada) es la única cuya violación
  sesga los coeficientes;
  \item la \textbf{homocedasticidad} y la \textbf{independencia} afectan a la varianza de
  $\hat{\boldsymbol\beta}$ y, por tanto, a los errores típicos, intervalos y contrastes;
  \item la \textbf{normalidad} sólo es necesaria para la inferencia exacta en muestras pequeñas;
  con $n$ grande, el teorema central del límite garantiza la normalidad aproximada de
  $\hat{\boldsymbol\beta}$.
\end{itemize}

\begin{table}[htbp]
\centering
\footnotesize
\caption{Diagnóstico de las hipótesis del modelo}\label{tab:diagnostico}
\begin{tabularx}{\linewidth}{l >{\raggedright\arraybackslash}X r r l}
\toprule
Hipótesis & Contraste & Estadístico & $p$ & Decisión al 5\,\% \\
\midrule
Linealidad & RESET de Ramsey (potencias 2 y 3) & \num{1.291} & \num{0.275} & no se rechaza $H_0$ \\
Normalidad & Shapiro-Wilk (residuos tipificados) & \num{0.9750} & $<\num{0.001}$ & se rechaza $H_0$ \\
 & Jarque-Bera & \num{952.9} & $<\num{0.001}$ & se rechaza $H_0$ \\
Homocedasticidad & Breusch-Pagan, MCO inicial & \num{262.7} & $<\num{0.001}$ & se rechaza $H_0$ \\
 & Breusch-Pagan, modelo final & \num{73.0} & $<\num{0.001}$ & se rechaza $H_0$ \\
 & Frente a 1/población, MCO inicial & \num{201.9} & $<\num{0.001}$ & se rechaza $H_0$ \\
 & Frente a 1/población, modelo final & \num{209.99} & $<\num{0.001}$ & se rechaza $H_0$ \\
Independencia & Durbin-Watson (orden del fichero) & \num{1.878} & --- & \\
 & Rachas de los signos & \num{-3.33} & $<\num{0.001}$ & se rechaza $H_0$ \\
 & Correlación intraclase por estado & \num{0.079} & $<\num{0.001}$ & se rechaza $H_0$ \\
Media cero & Media de los residuos & \num{-1.6e-13} & --- & por construcción \\
\bottomrule
\end{tabularx}
\par\smallskip{\footnotesize\raggedright La dependencia dentro de cada estado y la heterocedasticidad se absorben con errores típicos robustos por conglomerados; la falta de normalidad, con el bootstrap por conglomerados, que no la supone.\par}
\end{table}


\figura{diagnostico}{Gráficos de diagnóstico del modelo final: residuos tipificados frente a
valores ajustados, gráfico cuantil-cuantil normal, escala-localización y residuos estudentizados
frente al apalancamiento (el área de cada punto es proporcional a la distancia de Cook).}

\section{Linealidad}

El contraste RESET de \textcite{ramsey1969} añade al modelo las potencias segunda y tercera de los
valores ajustados: si la media condicional no fuera lineal en las explicativas, esos términos
captarían la curvatura omitida. No se rechaza la especificación ($F=\res{reset-f}$,
\res{reset-p}). Además, para cada explicativa continua se contrastó un término cuadrático con
corrección de Holm: el número de curvaturas significativas es \res{n-curvatura}
(\res{curvatura}). Los gráficos de residuos parciales (componente más residuo,
\cref{fig:residuos-parciales}) lo confirman: la curva suavizada sigue la recta del modelo en el
rango donde están casi todos los datos. Las transformaciones logarítmicas de la depuración
(D10) contribuyen a esta linealidad.

\figura{residuos-parciales}{Residuos parciales de las seis explicativas con mayor efecto: la recta
es la contribución estimada por el modelo y la curva, un suavizado LOWESS de los puntos.}

\section{Independencia de los errores}

El estadístico de \textcite{durbin1950},
\[
d=\frac{\sum_{i=2}^{n}(e_i-e_{i-1})^2}{\sum_{i=1}^{n}e_i^2}\approx 2(1-\hat\rho_1),
\]
vale \res{dw-final}, cerca de 2: no hay una autocorrelación de primer orden apreciable. Pero el
contraste se diseñó para series temporales, en las que el orden de las filas tiene sentido; en
datos transversales su resultado depende del orden arbitrario del fichero. Aquí el fichero está
agrupado por estado, de modo que la pregunta pertinente es si los condados de un mismo estado
tienen errores parecidos. Lo tienen:
\begin{itemize}
  \item la prueba de rachas sobre los signos de los residuos en el orden del fichero da $z=\res{rachas-res-z}$
  (\res{rachas-res-p}): menos rachas de las esperadas, es decir, residuos del mismo signo agrupados;
  \item la correlación intraclase de los residuos por estado es \res{icc-final}
  ($F=\res{icc-final-f}$, \res{icc-final-p}), y en un modelo mixto con intercepto aleatorio por estado
  la varianza entre estados representa una proporción \res{icc-mixto} de la varianza total.
\end{itemize}
La dependencia dentro del estado no sesga los coeficientes, pero sí hace que los errores típicos
clásicos sean demasiado pequeños: es la razón de usar errores típicos por conglomerados de estado
(\cref{sec:estructura-error}).

\section{Homocedasticidad}

El gráfico de residuos frente a valores ajustados no muestra un embudo marcado, pero los
contrastes detectan heterocedasticidad: Breusch-Pagan en la versión robusta de
\textcite{koenker1981} ($\mathrm{LM}=\res{bp-despues}$, \res{bp-despues-p}), White con los valores
ajustados y sus cuadrados ($\mathrm{LM}=\res{white-despues}$, \res{white-despues-p}) y, sobre
todo, el contraste frente a $1/n$ ($\mathrm{LM}=\res{bp-pob-despues}$, \res{bp-pob-despues-p}),
que identifica la fuente: la población del condado (\cref{sec:estructura-error}). El cociente de
Goldfeld-Quandt entre la varianza residual del tercio de condados más pequeños y la del tercio más
grande es \res{gq-despues}. La heterocedasticidad no sesga los coeficientes; la inferencia la
absorbe con los errores típicos robustos, y la estimación ponderada por la función de varianza
(\cref{cap:sensibilidad}) confirma que los coeficientes apenas cambian cuando se pondera.

\section{Normalidad de los residuos}

Los residuos tipificados tienen asimetría \res{res-asimetria} y curtosis en exceso
\res{res-curtosis}: son casi simétricos pero con colas más pesadas que la normal, como muestra el
gráfico Q-Q. Shapiro-Wilk ($W=\res{res-sw}$, \res{res-sw-p}) y Jarque-Bera
($\mathrm{JB}=\res{res-jb}$, \res{res-jb-p}) rechazan la normalidad exacta, como es esperable con
casi tres mil residuos y una mezcla de condados de tamaños muy distintos (los pequeños aportan los
residuos extremos). Las consecuencias son limitadas: los estimadores siguen siendo insesgados, y
con $n=\res{n-final}$ la distribución de $\hat{\boldsymbol\beta}$ es aproximadamente normal. Como
comprobación independiente, el bootstrap por conglomerados, que no supone normalidad, produce
intervalos muy parecidos a los paramétricos (\cref{cap:sensibilidad}).

\section{Multicolinealidad}

Tras la poda (\cref{sec:colinealidad}), el FIV máximo entre los efectos principales del modelo
final es \res{fiv-max-final} (tolerancia mínima por encima de 0,1) y el número de condición
\res{ncond-final}, lejos de los umbrales habituales (10 y 30). Las tolerancias y FIV de cada
coeficiente están en la \cref{tab:spss-coeficientes}. Los términos de interacción comparten
información con sus efectos principales por construcción, pero el centrado previo elimina la
colinealidad no esencial.

\section{Media de los residuos}

Al incluir la constante, los residuos MCO suman cero por construcción (su media es
\res{media-residuos}, error de redondeo). La hipótesis relevante,
$\E(\varepsilon\mid X)=0$, es la de linealidad, ya examinada.

\section{Observaciones atípicas e influyentes}\label{sec:influencia}

Se calcularon, para cada condado, el apalancamiento $h_{ii}$ (diagonal de la matriz sombrero
$H=X(X^{\top}X)^{-1}X^{\top}$), el residuo estudentizado externamente $r_i^{*}$, la distancia de
\textcite{cook1977},
\[
D_i=\frac{r_i^{2}}{k+1}\,\frac{h_{ii}}{1-h_{ii}},
\]
y los DFBETAS de cada coeficiente:
\begin{itemize}
  \item \res{n-palanca} condados tienen apalancamiento mayor que $2(k+1)/n=\res{umbral-palanca}$:
  combinaciones poco habituales de explicativas, no necesariamente problemáticas;
  \item \res{n-estudentizados} tienen $|r_i^{*}|>3$; el mayor es \res{mayor-residuo}
  ($r^{*}=\res{mayor-residuo-valor}$: mortalidad observada \res{mayor-residuo-y} frente a
  \res{mayor-residuo-ajuste} ajustada);
  \item \res{n-cook} superan el umbral orientativo $D_i>4/n=\res{umbral-cook}$, pero
  \res{n-cook-f50} superan la mediana de la $F$ de referencia, el criterio de influencia grave: el
  más influyente, \res{influyente1}, tiene $D=\res{influyente1-cook}$.
\end{itemize}

\begin{table}[htbp]
\centering
\footnotesize
\caption{Los diez condados más influyentes (distancia de Cook)}\label{tab:influyentes}
\begin{tabularx}{\linewidth}{>{\raggedright\arraybackslash}X r r r r r}
\toprule
Condado & $y$ & $\hat y$ & $r^*$ & $h_{ii}$ & Cook \\
\midrule
Williamsburg city, Virginia & \num{162.1} & \num{270.5} & \num{-6.05} & \num{0.051} & \num{0.0977} \\
De Baca County, New Mexico & \num{243.6} & \num{158.7} & \num{4.75} & \num{0.060} & \num{0.0716} \\
Crowley County, Colorado & \num{110.7} & \num{164.2} & \num{-3.02} & \num{0.082} & \num{0.0407} \\
Nome Census Area, Alaska & \num{277.6} & \num{201.4} & \num{4.19} & \num{0.027} & \num{0.0246} \\
Wibaux County, Montana & \num{214.4} & \num{141.1} & \num{4.03} & \num{0.028} & \num{0.0232} \\
Madison County, Mississippi & \num{292.5} & \num{158.4} & \num{7.35} & \num{0.008} & \num{0.0220} \\
Sioux County, North Dakota & \num{259.1} & \num{210.0} & \num{2.71} & \num{0.038} & \num{0.0146} \\
Bethel Census Area, Alaska & \num{218.8} & \num{168.2} & \num{2.79} & \num{0.036} & \num{0.0145} \\
North Slope Borough, Alaska & \num{256.9} & \num{175.6} & \num{4.44} & \num{0.013} & \num{0.0128} \\
Pitkin County, Colorado & \num{59.7} & \num{95.9} & \num{-2.01} & \num{0.054} & \num{0.0115} \\
\bottomrule
\end{tabularx}
\par\smallskip{\footnotesize\raggedright $r^*$: residuo estudentizado externamente; $h_{ii}$: apalancamiento.\par}
\end{table}


Ningún condado domina el ajuste. Aun así, se reestimó el modelo sin los \res{n-cook} condados con
$D_i>4/n$ y con regresión robusta de Huber, que pondera a la baja \res{n-huber} condados con
residuos grandes en lugar de eliminarlos (\cref{cap:sensibilidad}): las conclusiones no cambian.
Los atípicos detectados son condados reales, por lo que se mantienen en el modelo.
